\section{Related literature}

The research technology belongs to the tradition of endogenous and
semi-endogenous growth models in which resources devoted to ideas determine the
evolution of productivity \citep{romer1990,aghionhowitt1992,jones1995}. The
population term matters for the same reason it matters in \citet{jones1995}: when
human researchers remain essential, growth in the research workforce can anchor
long-run innovation. The present model asks when that anchor disappears because
capability-augmented research services can substitute for human researchers. Its
homothetic CES research technology separates the elasticity of
substitution from returns to research scale. This distinction preserves the
diminishing-returns margin in idea production emphasized by \citet{jones1995}
while allowing the model to study whether the input mix becomes increasingly
automated.

The paper is closely related to work on AI and economic growth.
\citet{aghionjonesjones2019} organize the possibility that automation of the idea
production function may generate very rapid growth and use ``singularity'' for a
finite-time divergence. \citet{nordhaus2021} develops empirical diagnostics for
an economic singularity, while \citet{trammellkorinek2023} studies growth under
transformative AI. \citet{jones2025rd} focuses directly on AI in research and
development. Our contribution is narrower but more explicit about equilibrium
prices and factor substitution: the same capability stock augments compute used
for inference and compute used to improve AI.

The terminology used to distinguish inference output, training activity,
algorithms, and model capital is informed by the economic accounting proposed by
\citet{korinekmckelvey2026}. Their framework treats compute as a central input
allocated between inference and training. In our notation, raw inference compute
and raw research compute are paid uses of final output; current AI capability
transforms them into economically useful services. The stock of capability is
closer to an index of model quality and algorithmic efficiency than to a physical
quantity of chips. The flow of capability improvement corresponds to the formation
of better model capital, broadly construed to include training, post-training,
experimentation, and algorithmic advances.

The final-production block connects to the older literature on the elasticity of
substitution and growth \citep{pitchford1960,jonesmanuelli1990,rebelo1991}. CES
production is not a cosmetic generalization of Cobb--Douglas: whether the
elasticity lies below or above one governs whether labor remains a bottleneck or
AI services can asymptotically dominate along the paths studied here. Empirical work such as
\citet{duffypapageorgiou2000} illustrates why the aggregate elasticity should not
be imposed without scrutiny. \citet{aumshin2024} sharpen this point by estimating
complementarity between equipment and labor but gross substitution between
software and labor. Their estimates do not calibrate our AI-service elasticity,
but they support separating software-like inputs from undifferentiated physical
capital.

The wage and factor-share results relate to the task and automation literature.
\citet{acemoglurestrepo2019} emphasize displacement and reinstatement of labor in
a task framework, and \citet{acemoglu2025} studies the macroeconomic gains from AI
through task exposure. That literature distinguishes tasks---units of work used
in production---from jobs, which bundle several tasks. Automation can therefore
change the tasks performed within a job without eliminating the job itself. Our
model abstracts from task content, occupations, and the creation of new human
tasks. Labor is an aggregate input and full employment holds by construction.
The results should therefore be interpreted as model implications for wages,
factor shares, and sectoral labor allocation, not as predictions about the number
of jobs displaced. Within that abstraction, the model shows how a wage can rise
even when labor income as a share of output falls because the scale of output
changes endogenously.
This mechanism is related to the distinction between wages and labor income shares
emphasized by \citet{autorkausik2026}, but here it is generated by AI research
feedback and monopoly supply. \citet{trammellkorinek2023} likewise show that wages
can rise or fall when output expands and labor income as a share of output contracts. Our conditional
AI-dominated limit maps that ambiguity into an explicit threshold for the
elasticity between labor and AI services.

Finally, the integrated developer connects the analysis to the emerging literature
on AI market structure \citep{korinekvipra2025,gans2024,demireretal2025}. Monopoly
and integration enter through separate channels: the developer charges a markup
on production AI services and internalizes the future operating profits created
by better capability. The present paper
holds market structure fixed in order to isolate the technology. Comparing this
benchmark with monopolistic competition is a natural next step.
